\documentclass[preprint,12pt,authoryear]{elsarticle}
\usepackage[margin=1in]{geometry}
\usepackage[T1]{fontenc}
\usepackage{lmodern,microtype}
\usepackage{amsmath,amssymb,amsthm}
\usepackage{booktabs,array,tabularx}
\usepackage{graphicx,flafter,placeins}
\biboptions{authoryear,round}
\usepackage[hidelinks]{hyperref}
\usepackage{caption}
\captionsetup{font=small,labelfont=bf}
\hypersetup{pdftitle={Training-Derived Regions of Nearest-Prototype Agreement in Gaussian RBF Classifiers},pdfauthor={Meshaal Mouawad}}
\newtheorem{theorem}{Theorem}
\newtheorem{lemma}{Lemma}
\newtheorem{proposition}{Proposition}
\newtheorem{corollary}{Corollary}
\newcommand{\R}{\mathbb{R}}
\newcommand{\norm}[1]{\left\lVert#1\right\rVert}
\newcommand{\tr}{\mathsf{T}}
\newcommand{\sharpb}{\mathrm{sharp}}
\newcommand{\basicb}{\mathrm{basic}}
\newcommand{\frobb}{\mathrm{F}}
\widowpenalty=10000
\clubpenalty=10000
\displaywidowpenalty=10000
\setlength{\parskip}{0.2em}
\setlength{\emergencystretch}{2em}
\journal{Neural Networks}
\begin{document}
\begin{frontmatter}
\title{Training-Derived Regions of Nearest-Prototype Agreement in Gaussian RBF Classifiers}
\author[msu]{Meshaal Mouawad\corref{cor1}}
\ead{mm4922@msstate.edu}
\cortext[cor1]{Corresponding author.}
\affiliation[msu]{organization={Department of Electrical and Computer Engineering, Mississippi State University},city={Mississippi State},postcode={39762},state={MS},country={USA}}
\begin{abstract}
We establish training-derived conditions that certify or exclude nearest-prototype agreement in Gaussian radial basis function classifiers over continuous bandwidth--penalty regions. A class-wise residual decomposition, prototype-block spectral floor, and exact penalty envelope certify complete regions $(0,t]\times[\lambda_0,\infty)$. A strict competing vote in the large-penalty limit supplies a finite disagreement bound, and a signed interval construction certifies local bandwidth bands with finite penalty floors. A frozen MNIST and Fashion-MNIST study spans five seeds and 50--200 prototypes. Immediately beyond the sufficient frontier, more than 99.5\% of query--region pairs remain unresolved. At the median-distance bandwidth scale, the original witness union excludes 88.10\% and 84.48\% of queries, respectively. Local bands at half this scale certify 16.88\% and 40.09\% at ten times the endpoint Gram-matrix norm. Matched-coverage class analysis reveals concentrated predictions within some certified subsets. Across all 30 dataset--seed--budget groups, 328 readouts retain at least 95\% of their group's best validation accuracy; all have zero full-width certificate coverage. Exhaustive reconstruction of 3,366,480 query--width paths identifies 190,405 returns to agreement after sampled disagreement. Expanded finite probes recover at least 95\% of the original limiting-witness cases after mean 15.6 and 20.2 checks at the median-distance scale. These results establish computable continuous decision-persistence guarantees and characterize their empirical separation from competitive prediction, finite-path behavior, and selected-set correctness.
\end{abstract}
\begin{keyword}
radial basis function classifier \sep nearest prototype \sep Gaussian bandwidth \sep ridge regression \sep uniform certificate \sep parameter region
\end{keyword}
\end{frontmatter}
\section{Introduction}
The label of the nearest prototype does not always describe a Gaussian radial basis function (RBF) classifier's decision. Its learned output layer combines responses from all prototypes, with coefficients determined by the full labeled training sample. We establish training-derived conditions that determine when the nearest-prototype label persists throughout a continuous bandwidth--penalty region, when that persistence is impossible, and which queries remain unresolved. The proofs give both sufficient agreement conditions and finite disagreement thresholds. Experiments on MNIST and Fashion-MNIST show that sufficient certification disappears before widespread identified failures of the requested family, local width bands admit additional guarantees at large penalties, and all 328 archived readouts meeting the 95\% validation-competitive rule have zero full-width coverage. The analysis therefore locates both the guaranteed persistence regions and their empirical separation from high-accuracy operating points.

A prototype is a reference input used to measure similarity. Here, prototypes are selected training images. Each Gaussian response decreases with distance from its prototype, and its bandwidth controls how broadly that response extends. Ridge regression learns the coefficients that combine these responses into class scores. The separation between nonlinear basis functions and linear output estimation is central to the RBF formulation of Broomhead and Lowe \citeyearpar{broomhead1988}. It also creates the present decision-summary problem: the largest individual response need not identify the class with the largest weighted score.

Existing research establishes several relevant mechanisms. Coefficient-regularized RBF networks have a data-dependent kernel interpretation and statistical error bounds \citep{que2016}. Gaussian narrow-bandwidth limits have a nearest-neighbour dominance mechanism, also discussed in neural-network kernel analyses \citep{radhakrishnan2023}. Prototype architectures can impose a relation between components and predicted classes \citep{saralajew2025}, and learned Gaussian geometry can reveal feature importance \citep{dagostino2024}. Regularization-path methods bound optimization or validation behavior over parameter intervals \citep{shibagaki2015,ndiaye2019}. Our question concerns a fixed conventional ridge RBF family: which computable training conditions establish or rule out equality with a particular nearest-prototype label as both bandwidth and regularization vary?

The parameter region specifies what the summary must survive. For a fixed query, the first region includes every positive bandwidth up to an upper endpoint and every finite ridge penalty above a floor. This asks whether proximity remains a valid label summary as responses broaden and regularization changes the readout. Large penalties expose a specific structural test: after positive rescaling, the class scores converge to training-label votes assembled through prototype features. A competing class that strictly wins this vote at any included width forces a finite disagreement inside the requested region. Thus an unbounded penalty interval supplies a way to distinguish failure of the family property from failure of a sufficient inequality.

We first compare the fitted weights with weights that assign each prototype to its own class. An exact residual decomposition separates prototype overlap, non-prototype labeled responses, and a quadratic overlap correction. A spectral floor and a two-endpoint envelope then bound the effect of that residual throughout $(0,t]\times[\lambda_0,\infty)$. Comparing the bound with a query margin certifies agreement without fitting a readout. The large-penalty test complements this sufficient condition with an exclusion criterion and an explicit finite witness threshold. Their outputs classify each query and requested region as certified, excluded, or unresolved.

The limiting votes also support a second positive construction. A signed interval enclosure controls their finite-penalty correction over a local bandwidth band $[s,t]$, producing a certificate for $[s,t]\times[\Lambda,\infty)$ with a finite floor $\Lambda$. Together, the results provide three analytical capabilities: full-width sufficient certification with ordered class-wise and norm comparators; strict exclusion of impossible full-width regions; and local-width penalty-tail certification. The frozen image study measures the first capability and its bandwidth frontier. Focused post hoc analyses measure the exclusion partition and local certificate on the same 100-prototype groups, reporting region geometry, coverage, fitted agreement, and class accuracy separately. Follow-up comparisons locate sufficient reversal bounds relative to validation-based competitive models, measure exact overlap with all archived finite witnesses, and compare certified subsets with informed selectors at identical coverage. Class-balanced selection diagnostics identify which queries and classes account for certified-subset accuracy. Exhaustive reconstruction of all 30 original groups then measures observed penalty-path reversals, validation-competitive coverage, and stability across prototype budgets and seeds. A fixed-order replay identifies how finite exclusions accumulate relative to limiting witnesses. Figure~\ref{fig:paths} introduces the two decisions whose equality is studied.

\begin{figure}[!htbp]
\centering
\includegraphics[width=\textwidth]{Figure_01_Two_Decision_Paths.pdf}
\caption{Nearest-prototype reference and learned Gaussian radial basis function (RBF) prediction. (A) Squared distances $d_j^2=\norm{x-c_j}_2^2$ from query $x$ to selected prototypes $c_j$ identify a nearest index $j_0$ and its reference class $a$. (B) Training inputs $X$ and labels $Y$ supply Gaussian training features $H_\sigma$ and a ridge readout $W_{\sigma,\lambda}$; query features $h_\sigma(x)$ produce the fitted class $\hat y_{\sigma,\lambda}$. The question mark compares the two output labels. Theorem~\ref{thm:agreement} gives a sufficient condition for their equality throughout a bandwidth--penalty region. This diagram explains the two decision paths.}
\label{fig:paths}
\end{figure}

\section{Related work}
\subsection{Prototype meaning and fitted class decisions}
Prototype networks make a reference example or learned representative part of the prediction mechanism. The relation between that representative and the final class depends on the output rule. Saralajew and colleagues \citeyearpar{saralajew2025} examine this relation for RBF classifiers and show why unconstrained coefficients can separate the nearest components from the predicted class. Their response constructs a probabilistic network with an interpretable prototype-class relationship and analyzes robustness to input perturbations. The present study keeps a conventional ridge readout and asks when the fitted decision has a particular prototype-label summary.

Bugmann \citeyearpar{bugmann1998} relates normalized Gaussian RBF responses to a Voronoi-like partition when one hidden node dominates and illustrates nearest-center behavior in classification. For fixed coefficients, normalizing all class scores by the same positive total response preserves the winning class. The question pursued here concerns when a ridge-fitted readout shares the nearest selected training input's label throughout a bandwidth--penalty region.

Nearest-neighbour limiting behaviour also appears in kernel-based analyses of neural networks. Radhakrishnan, Belkin, and Uhler \citeyearpar{radhakrishnan2023} prove a 1-NN regime for infinitely wide and deep networks and explicitly relate their kernel-ratio argument to the Gaussian narrow-bandwidth limit. This identifies a known limiting mechanism. The present question concerns computable finite-bandwidth regions for a coefficient-ridge RBF readout, with agreement required throughout a specified penalty interval.

Fazaeli-Javan and Ghiasi-Shirazi \citeyearpar{fazaeli2026} investigate how training samples contribute to softmax-neuron decisions through learned prototypes. Their positive and negative distance-based representations connect a neuron's decision to prototype influences and supporting or opposing samples. They report that a small set of near-boundary samples disproportionately shapes the prototypes and decision boundaries. This supplies a sample-level account of learned coefficients. Our analytical target is the equality of two labels across specified values of bandwidth and regularization: the label returned by a fixed ridge RBF family and the label attached to its nearest selected input. Training labels enter the analysis through class-specific response sums, so examples that are not selected as prototypes remain part of the certificate.

These studies establish why prototype geometry alone is insufficient to characterize a general learned readout. They also identify useful distinctions between explanatory targets. A decomposition of sample influence describes how a score is assembled; a constrained architecture defines how prototypes determine a decision; a parameter-region certificate establishes when a conventional fitted rule shares a selected prototype's label. The last operation is the question addressed here.

\subsection{Gaussian scale and numerical structure}
Que and Belkin \citeyearpar{que2016} analyze coefficient-norm regularization of RBF networks and interpret the resulting models through data-dependent kernels. They give statistical error bounds for training-point centers and weighted k-means centers, with quantization error entering the latter analysis. This establishes a theoretical basis for the regularized readout studied here. Our target is a query's agreement with its nearest labeled training prototype over a continuous bandwidth--penalty region, established through the fixed training matrix and its selected-row structure.

D'Agostino, Ilievski, and Shoemaker \citeyearpar{dagostino2024} equip Gaussian RBF networks with a learnable precision matrix and extract feature-importance and active-subspace information from its spectrum. Their analysis concerns learned representation geometry. Fixed prototype coordinates and a shared scalar width expose a different analytical question here: how training geometry determines sufficient agreement conditions and large-penalty obstructions for a fitted class decision.

Gaussian scale controls overlap between basis functions and changes the matrix used to estimate output weights. Wurzberger and Schwenker \citeyearpar{wurzberger2024} study initialization and learning in deep RBF networks, including clustering, covariance estimation, and partially connected distance calculations. Their work places the geometry of the representation within the learning procedure. Fixing prototype inputs and using one shared width exposes a simpler structure: every off-diagonal response increases as width increases, and the prototype-to-prototype block remains symmetric.

Noorizadegan, Wang, and Ling \citeyearpar{noorizadegan2026reliable} analyze Gaussian Kolmogorov--Arnold networks through the conditioning of Gaussian feature blocks and the choice of scale. Their study links reliable computation to the relation between scale and grid geometry in function approximation and differential-equation tasks. Noorizadegan's \citeyearpar{noorizadegan2026partition} partition-of-unity construction normalizes Gaussian edge features, reproduces constants, and studies scale sensitivity. Together, these papers make scale a mathematical property of the representation rather than an isolated tuning number. They motivate examining what Gaussian overlap permits one to guarantee about decisions.

Our use of scale concerns a fixed set of labeled training prototypes. Symmetry yields an exact cancellation in the readout residual, and increasing overlap supplies endpoint bounds for every smaller width. The certificate combines this geometric information with training-label contributions. Its output comprises sufficient agreement regions and strict obstructions to agreement for a specific query.

\subsection{The object varied by a prediction certificate}
A prediction certificate must specify what is allowed to change. Voracek and Hein \citeyearpar{voracek2022} study adversarial robustness of nearest-prototype classifiers and derive guarantees concerning input perturbations. Rosenfeld and colleagues \citeyearpar{rosenfeld2020} certify predictions against changes to training labels using randomized smoothing. In both cases, the varying object defines the statement proved. Here, query coordinates, training coordinates, prototype identities, and labels remain fixed; the Gaussian bandwidth and ridge penalty vary.

Regularization-path analysis already supplies guarantees over parameter intervals. Shibagaki and colleagues \citeyearpar{shibagaki2015} derive validation-score bounds from an approximate solution at another regularization value and use them to bound cross-validation error over an interval for regularized binary classification. Ndiaye and colleagues \citeyearpar{ndiaye2019} construct approximate solution paths and use duality-gap safe regions to bound validation error along a regularization path with a fixed design matrix. Their results establish parameter-range reasoning for model selection. The construction here varies the Gaussian feature matrix with bandwidth as well as varying the ridge penalty, and uses an explicit prototype-label reference to certify nearest-prototype agreement over the resulting two-parameter region.

The mathematical tools include residual-based solution bounds, norm inequalities, and eigenvalue estimates. For a nonsingular linear system, multiplying the residual by the inverse matrix expresses the solution error; Bindel \citeyearpar{bindel2025} presents this standard residual-based argument. The present derivation uses the selected-row structure of a Gaussian feature matrix to make the residual nonnegative and label-specific after majorization. It then maximizes the relevant bound over the entire penalty interval. The same training features supply limiting class votes, a finite reversal threshold, and signed endpoint enclosures for local-width agreement. Table~\ref{tab:prior} summarizes the resulting distinction in the scientific object and output.

\begin{table}[!htbp]
\centering\small
\caption{Closest research relationships. The comparison identifies the object studied and the result supplied by each approach.}
\label{tab:prior}
\begin{tabularx}{\textwidth}{@{}>{\raggedright\arraybackslash}p{0.23\textwidth}>{\raggedright\arraybackslash}X>{\raggedright\arraybackslash}X@{}}
\toprule
Work & Object studied & Analytical or explanatory output \\
\midrule
\citet{bugmann1998} & Normalized Gaussian RBFs with learned centers & Voronoi-like response partition and classification examples \\
\citet{radhakrishnan2023} & Infinite-width/depth neural-network kernels & 1-NN regime and Gaussian ratio-dominance analogy \\
\citet{dagostino2024} & Gaussian RBFs with learned precision matrices & Active-subspace and feature-importance analysis \\
\citet{que2016} & Coefficient-regularized RBFs and k-means centers & Data-dependent kernels and statistical error bounds \\
\citet{saralajew2025} & Probabilistic prototype architecture and perturbed inputs & Prototype interpretation and input-robustness analysis \\
\citet{fazaeli2026} & Learned prototypes in softmax neurons & Supporting and opposing sample influences \\
\citet{noorizadegan2026reliable} & Gaussian scale and feature-block geometry & Conditioning and reliability analysis \\
\citet{noorizadegan2026partition} & Normalized Gaussian edge features & Partition-of-unity representation and scale-sensitivity analysis \\
\citet{voracek2022} & Perturbed inputs to nearest-prototype classifiers & Adversarial prediction guarantees \\
\citet{rosenfeld2020} & Changed training labels & Label-flip prediction certificates \\
\citet{shibagaki2015} & Regularized binary classification with fixed inputs & Validation-score and cross-validation-error bounds over a regularization interval \\
\citet{ndiaye2019} & Regularization path with a fixed design matrix & Solution-path approximation and validation-error bounds \\
Present method & Width and ridge penalty with fixed data and prototypes & Sufficient agreement regions, strict exclusions, and local penalty-tail certificates \\
\bottomrule
\end{tabularx}
\end{table}

\section{Model and training-derived quantities}
\subsection{Purpose and model definition}
The analysis takes a training set, selected prototypes, a query, and a requested bandwidth--penalty region. Sufficient conditions certify that the entire region predicts the query's nearest-prototype label; strict disagreement tests exclude that property. Cases receiving neither conclusion remain unresolved. The calculation first summarizes overlap among labeled training inputs at the upper bandwidth. It then compares a bound derived from these summaries with the query's separation from competing-class prototypes.

Let $x_i\in\R^d$, $i=1,\ldots,N$, be fixed training inputs, and let $Y\in\{0,1\}^{N\times C}$ contain one-hot labels for $C$ classes. Select $K$ distinct training-row indices as prototypes $c_j$. The selector $S\in\{0,1\}^{N\times K}$ places a one at each selected row, so $S^\tr S=I_K$, and $Y_C=S^\tr Y$ contains prototype labels. Selection concerns distinct rows; different rows may have equal coordinates. For a shared positive bandwidth $\sigma$, the feature matrix is
\begin{equation}
 H_\sigma(i,j)=\exp\left[-\frac{\norm{x_i-c_j}_2^2}{2\sigma^2}\right].
 \label{eq:features}
\end{equation}
The readout minimizes the sum-of-squares objective
\begin{equation}
 \norm{H_\sigma W-Y}_F^2+\lambda\norm{W}_F^2,\qquad \lambda\geq0,
 \label{eq:objective}
\end{equation}
with no intercept. Write $G_{\sigma,\lambda}=H_\sigma^\tr H_\sigma+\lambda I_K$. When $G_{\sigma,\lambda}$ is positive definite, the solution is
\begin{equation}
 W_{\sigma,\lambda}=G_{\sigma,\lambda}^{-1}H_\sigma^\tr Y.
 \label{eq:readout}
\end{equation}
For a query $x$, the feature column $h_\sigma(x)$ is defined by the same Gaussian responses, and the largest entry of $h_\sigma(x)^\tr W_{\sigma,\lambda}$ determines its class. The certificate below supplies a positive eigenvalue bound on its accepted regions, including any accepted region with $\lambda_0=0$.

\subsection{Separating self-response from overlap}
Each selected prototype has response one at its own training row. Remove these self-responses by defining $E_\sigma=H_\sigma-S$. This matrix is entrywise nonnegative, and every entry is nondecreasing in $\sigma$. Its selected-row block $E_{C,\sigma}=S^\tr E_\sigma$ is symmetric with zero diagonal because all units use the same Euclidean Gaussian width. Let $R$ denote the remaining training rows. At the endpoint $t$, define
\begin{equation}
 V_t=E_{C,t}Y_C,\qquad
 P_t=E_{R,t}^\tr Y_R,\qquad
 Q_t=E_t^\tr(E_tY_C).
 \label{eq:terms}
\end{equation}
All three are $K\times C$ nonnegative matrices. The matrix $V_t$ groups overlap between prototypes by class; $P_t$ groups responses from the other labeled observations; $Q_t$ records a two-stage overlap correction relative to the prototype-label reference. Their entries increase with $t$, so endpoint values bound the corresponding quantities at every smaller width. The parenthesization of $Q_t$ allows its computation without forming $E_t^\tr E_t$.

The prototype block also supplies a spectral floor. Let
\begin{equation}
 B_t=S^\tr H_t=I_K+E_{C,t},\qquad
 \beta_t=\max_j\sum_{k\ne j}B_t(j,k),\qquad
 \gamma_t=\bigl[\max(1-\beta_t,0)\bigr]^2.
 \label{eq:spectral}
\end{equation}
Here $\beta_t$ is the largest total off-diagonal prototype response, and $0\leq\gamma_t\leq1$. The clipping precedes the square. Strong overlap may reduce this floor to zero; a positive ridge penalty still supplies a positive denominator for the subsequent bound.

\subsection{An exact readout residual}
The reference weights $W_0(\lambda)=Y_C/(1+\lambda)$ place a coefficient $1/(1+\lambda)$ on each prototype's own class and zero on its other classes. They provide a tractable comparison with the learned readout. Their normal-equation residual is
\begin{equation}
 \mathcal R_{\sigma,\lambda}
 =H_\sigma^\tr Y-G_{\sigma,\lambda}W_0(\lambda).
 \label{eq:residualdef}
\end{equation}
The following identity separates the three training contributions before any norm bound is applied.

\begin{lemma}[Residual decomposition]
For the model in Eqs.~\eqref{eq:features}--\eqref{eq:objective},
\begin{equation}
 \mathcal R_{\sigma,\lambda}
 =\frac{\lambda-1}{\lambda+1}V_\sigma
   +P_\sigma-\frac{Q_\sigma}{\lambda+1}.
 \label{eq:residual}
\end{equation}
When $G_{\sigma,\lambda}$ is invertible, the difference $\Delta W=W_{\sigma,\lambda}-W_0(\lambda)$ equals $G_{\sigma,\lambda}^{-1}\mathcal R_{\sigma,\lambda}$.
\end{lemma}
\begin{proof}
The decomposition $H_\sigma=S+E_\sigma$ gives $H_\sigma^\tr Y=Y_C+V_\sigma+P_\sigma$. Symmetry of $E_{C,\sigma}$ gives
\[
 H_\sigma^\tr H_\sigma
 =I_K+2E_{C,\sigma}+E_\sigma^\tr E_\sigma.
\]
Substitution into Eq.~\eqref{eq:residualdef} cancels $Y_C$ and yields Eq.~\eqref{eq:residual}. Since $G_{\sigma,\lambda}W_{\sigma,\lambda}=H_\sigma^\tr Y$, the residual definition gives $G_{\sigma,\lambda}(W_{\sigma,\lambda}-W_0)=\mathcal R_{\sigma,\lambda}$, proving the second statement.
\end{proof}

Equation~\eqref{eq:residual} retains a cancellation associated with the symmetric prototype block. In particular, the coefficient of $V_\sigma$ vanishes at $\lambda=1$. The next section uses its absolute value to bound class-pair changes while preserving the distinction between prototype and non-prototype contributions.

\section{Uniform nearest-prototype agreement}
\subsection{A bound for the entire parameter region}
For different classes $a$ and $b$, write $V_{ab,t}=V_t[:,a]+V_t[:,b]$, and define $P_{ab,t}$ and $Q_{ab,t}$ in the same way. These nonnegative vectors bound the absolute difference between the corresponding class contributions. The first step controls the inverse matrix in the residual identity throughout the region.

\begin{lemma}[Spectral floor]
For every $0<\sigma\leq t$ and $\lambda\geq0$,
\begin{equation}
 \lambda_{\min}(G_{\sigma,\lambda})\geq\gamma_t+\lambda.
 \label{eq:eigenlower}
\end{equation}
\end{lemma}
\begin{proof}
The selected rows form $B_\sigma$, so $H_\sigma^\tr H_\sigma\succeq B_\sigma^\tr B_\sigma$. If $\beta_t<1$, symmetry and the Gershgorin bound imply $\lambda_{\min}(B_\sigma)\geq1-\beta_\sigma\geq1-\beta_t>0$. Squaring this bound and adding $\lambda I_K$ proves Eq.~\eqref{eq:eigenlower}. If $\beta_t\geq1$, then $\gamma_t=0$ and positive semidefiniteness of $H_\sigma^\tr H_\sigma$ gives the stated ridge floor.
\end{proof}

Assume $\gamma_t+\lambda_0>0$. The residual identity, entrywise monotonicity, and Eq.~\eqref{eq:eigenlower} imply
\begin{equation}
 (1+\lambda)\norm{\Delta W[:,a]-\Delta W[:,b]}_2
 \leq\norm{F_{ab}(\lambda;t)}_2,
 \label{eq:pairperturb}
\end{equation}
where the nonnegative vector
\begin{equation}
 F_{ab}(\lambda;t)=
 \frac{|\lambda-1|V_{ab,t}+(\lambda+1)P_{ab,t}+Q_{ab,t}}
 {\gamma_t+\lambda}
 \label{eq:majorant}
\end{equation}
uses training information only at $t$. Thus width has already been controlled; the remaining task is to bound the norm of this vector for every penalty above the floor.

\begin{proposition}[Two-endpoint penalty envelope]
For $t>0$, $\lambda_0\geq0$, and $\gamma_t+\lambda_0>0$, define
\begin{equation}
 D^{\sharpb}_{ab}(t,\lambda_0)=
 \max\left\{\norm{F_{ab}(\lambda_0;t)}_2,
                  \norm{V_{ab,t}+P_{ab,t}}_2\right\}.
 \label{eq:sharp}
\end{equation}
Then $D^{\sharpb}_{ab}$ is the exact supremum of $\norm{F_{ab}(\lambda;t)}_2$ over $\lambda\geq\lambda_0$, and Eq.~\eqref{eq:pairperturb} holds with $D^{\sharpb}_{ab}$ on its right throughout $(0,t]\times[\lambda_0,\infty)$.
\end{proposition}
\begin{proof}
For $\lambda\geq\lambda_0$, the scalar triangle inequality gives
$|\lambda-1|\leq|\lambda_0-1|+\lambda-\lambda_0$. Set
$\theta=(\gamma_t+\lambda_0)/(\gamma_t+\lambda)\in(0,1]$. Entrywise,
\[
 F_{ab}(\lambda;t)
 \leq\theta F_{ab}(\lambda_0;t)
       +(1-\theta)(V_{ab,t}+P_{ab,t}).
\]
The Euclidean norm is monotone for nonnegative vectors and convex, so its value is bounded by the larger endpoint norm in Eq.~\eqref{eq:sharp}. The first norm occurs at $\lambda_0$ and the second is approached as $\lambda\to\infty$. Their maximum is therefore the supremum of this majorant. Equation~\eqref{eq:pairperturb} transfers the bound to the class-pair weight difference.
\end{proof}

The exact maximization in this proposition concerns the chosen nonnegative majorant. It removes a continuous penalty search from evaluating that bound. The nearest-prototype theorem now follows by comparing the resulting weight perturbation with a query margin.

\subsection{From a class-pair bound to a label guarantee}
Let $j_0$ index a nearest prototype to a query $x$, let its label be $a$, and let $d_j^2=\norm{x-c_j}_2^2$. Normalize all query responses by the positive nearest response. At the upper bandwidth, this produces
\begin{equation}
 r_j(t)=\exp\left[-\frac{d_j^2-d_{j_0}^2}{2t^2}\right],\qquad
 n_t=\norm{r(t)}_2,\qquad
 L_b(t)=1-\sum_{j:Y_C(j,b)=1}r_j(t).
 \label{eq:query}
\end{equation}
The nearest response equals one. The quantity $L_b(t)$ subtracts the total response of competing-class prototypes and lower-bounds the reference margin against class $b$. It is computed from prototype distances and labels without evaluating a fitted readout.

\begin{theorem}[Agreement on a complete parameter region]
Suppose $t>0$, $\lambda_0\geq0$, and $\gamma_t+\lambda_0>0$. If
\begin{equation}
 L_b(t)>n_tD^{\sharpb}_{ab}(t,\lambda_0)
 \quad\text{for every }b\ne a,
 \label{eq:certificate}
\end{equation}
then, for every $0<\sigma\leq t$ and every finite $\lambda\geq\lambda_0$, the unique winning class of $h_\sigma(x)^\tr W_{\sigma,\lambda}$ is $a$.
\label{thm:agreement}
\end{theorem}
\begin{proof}
Dividing $h_\sigma(x)$ by its nearest response multiplies all class scores by the same positive scalar and produces $r(\sigma)$. Squared-distance gaps are nonnegative, so $r_j(\sigma)\leq r_j(t)$ when $\sigma\leq t$. The prototype-reference margin satisfies
\begin{align*}
 (1+\lambda)r(\sigma)^\tr\bigl(W_0[:,a]-W_0[:,b]\bigr)
 &=r(\sigma)^\tr\bigl(Y_C[:,a]-Y_C[:,b]\bigr)\\
 &\geq 1-\sum_{j:Y_C(j,b)=1}r_j(t)=L_b(t).
\end{align*}
The inequality retains the response of prototype $j_0$ and drops the other nonnegative responses from class $a$. By Cauchy--Schwarz and the penalty envelope,
\[
 (1+\lambda)\left|r(\sigma)^\tr
       \bigl(\Delta W[:,a]-\Delta W[:,b]\bigr)\right|
 \leq \norm{r(\sigma)}_2D^{\sharpb}_{ab}
 \leq n_tD^{\sharpb}_{ab}.
\]
The strict condition in Eq.~\eqref{eq:certificate} consequently leaves a positive learned margin against every competing class throughout the region. Positive query scaling preserves this unique winning class in the original scores.
\end{proof}

Figure~\ref{fig:mechanism} connects the training-only bound to the query margin and shows the region guaranteed when every class comparison passes. The procedure abstains if the spectral denominator is zero or any strict inequality fails. A nearest tie across classes forces a competing margin to be nonpositive and therefore prevents certification. A class absent from the prototype set remains a competitor: its training labels still contribute through $P_t$. Thus the construction accounts for the full training-label matrix, including classes represented only outside the selected rows.


\begin{figure}[!htbp]
\centering
\includegraphics[width=\textwidth]{Figure_02_Uniform_Certificate_Mechanism.pdf}
\caption{Uniform nearest-prototype agreement certificate. (A) Prototype overlap $V_t$, non-prototype labeled responses $P_t$, quadratic overlap $Q_t$ in Eq.~\eqref{eq:terms}, and spectral floor $\gamma_t$ in Eq.~\eqref{eq:spectral} determine the class-pair bound $D^{\sharpb}_{ab}(t,\lambda_0)$ in Eq.~\eqref{eq:sharp}. It requires $t>0$, $\lambda_0\geq0$, and $\gamma_t+\lambda_0>0$. Query distances determine the reference class $a$, competing-class margin $L_b(t)$, and response norm $n_t$ in Eq.~\eqref{eq:query}. The strict comparison must pass for every $b\ne a$. (B) A passed test guarantees class $a$ for every $0<\sigma\leq t$ and finite $\lambda\geq\lambda_0$. The shaded region illustrates this mathematical domain; empirical coverage is measured in Fig.~\ref{fig:coverage}.}
\label{fig:mechanism}
\end{figure}

\subsection{Comparators and nested certified regions}
Two bounds with the same region semantics isolate the role of the retained class-wise information and the two-endpoint calculation. Replacing $|\lambda-1|$ by $1+\lambda$ yields
\begin{equation}
 D^{\basicb}_{ab}=
 \norm{\frac{(1+\lambda_0)(V_{ab,t}+P_{ab,t})+Q_{ab,t}}
 {\gamma_t+\lambda_0}}_2.
 \label{eq:basic}
\end{equation}
For a norm/count comparator, let $\eta_t=\norm{E_t}_F$, let $n_{ab}$ count training labels in $\{a,b\}$, and let $k_{ab}$ count prototype labels in those classes. The cancellation-aware Frobenius bound is
\begin{equation}
 D^{\frobb}_{ab}=
 \frac{(1+\lambda_0)\eta_t\sqrt{n_{ab}}+\eta_t^2\sqrt{k_{ab}}}
 {\gamma_t+\lambda_0}.
 \label{eq:frobenius}
\end{equation}
Each comparator uses Eq.~\eqref{eq:certificate} with its corresponding bound.

\begin{proposition}[Bound and certificate hierarchy]
For every applicable endpoint and class pair,
\begin{equation}
 D^{\sharpb}_{ab}\leq D^{\basicb}_{ab}\leq D^{\frobb}_{ab}.
 \label{eq:hierarchy}
\end{equation}
At $\lambda_0=0$, the sharp and basic bounds coincide. The queries certified by the Frobenius condition are contained in those certified by the basic condition, which are contained in those certified by the sharp condition.
\end{proposition}
\begin{proof}
The vector in Eq.~\eqref{eq:basic} dominates $F_{ab}(\lambda_0;t)$ because $|\lambda_0-1|\leq1+\lambda_0$. It also dominates $V_{ab,t}+P_{ab,t}$ because $\gamma_t\leq1$. Monotonicity of the norm proves the first inequality, including equality at zero penalty floor. Let $z=Y[:,a]+Y[:,b]$ and $z_C=Y_C[:,a]+Y_C[:,b]$. Symmetry of the prototype block gives $V_{ab,t}+P_{ab,t}=E_t^\tr z$, and $Q_{ab,t}=E_t^\tr E_tz_C$. Therefore
\[
 \norm{V_{ab,t}+P_{ab,t}}_2\leq\eta_t\sqrt{n_{ab}},
 \qquad
 \norm{Q_{ab,t}}_2\leq\eta_t^2\sqrt{k_{ab}}.
\]
The triangle inequality applied to Eq.~\eqref{eq:basic} proves the second inequality. Substitution into the same strict margin test proves inclusion of the certified sets.
\end{proof}

For a fixed penalty floor, every numerator in the bounds increases with $t$ and $\gamma_t$ decreases. The query norm and competing-class response sums also increase. Hence certification at $t$ implies certification at every smaller upper bandwidth. Raising the penalty floor shrinks the interval defining the exact supremum in Eq.~\eqref{eq:sharp}, so the sharp constant cannot increase. These nested regions permit a search for a certified bandwidth frontier using training quantities. The original endpoint grid measures this behavior at predefined widths; the supplementary analysis uses the same nesting to bracket each query's frontier.

After Gaussian feature construction, forming $V_t$, $P_t$, and $Q_t$ costs $O(NKC)$ arithmetic; preparing all class-pair constants costs $O(KC^2)$. Computing $Q_t$ as $E_t^\tr(E_tY_C)$ avoids the $O(NK^2)$ full training Gram construction. Row-wise accumulation can retain $O(KC+K^2)$ state, excluding inputs and query batches. Once query distances are available, class response sums require $O(K)$ work followed by $C-1$ comparisons. These counts describe the certificate calculation; the experiment retains the full feature matrix to fit comparison models.

\subsection{A necessary condition from the large-penalty limit}
\label{sec:ridge-obstruction}
An agreement certificate gives a sufficient condition for a complete parameter
region. The ridge limit supplies a complementary test that can rule such a
region out. For a fixed width, write
\[
 A_\sigma=H_\sigma^\tr H_\sigma,\qquad
 Z_\sigma=H_\sigma^\tr Y=Y_C+V_\sigma+P_\sigma,
 \qquad k_c(\sigma,x)=r(\sigma)^\tr Z_\sigma[:,c].
\]
Here $r(\sigma)$ is the nearest-response normalization in
Eq.~\eqref{eq:query}. The score $k_c$ sums training-label contributions
through the prototype features. It can be calculated from the same training
quantities used by the original certificate, without fitting output weights.

\begin{proposition}[Large-penalty decision and exclusion]
\label{prop:ridge-limit}
For every fixed positive width and query,
\[
 \lim_{\lambda\to\infty}\lambda r(\sigma)^\tr W_{\sigma,\lambda}
 =r(\sigma)^\tr Z_\sigma.
\]
Let $a$ be the nearest-prototype label. If, for some competing class $b$,
$\delta=k_b(\sigma,x)-k_a(\sigma,x)>0$, then every region
$(0,t]\times[\lambda_0,\infty)$ with $t\geq\sigma$ fails to preserve
label $a$ at this query. More explicitly, put
$v=Z_\sigma[:,b]-Z_\sigma[:,a]$. Class $b$ has a strictly larger fitted
score than class $a$ whenever
\[
 \lambda>\frac{\norm{r(\sigma)}_2\norm{A_\sigma v}_2}{\delta}.
\]
\end{proposition}
\begin{proof}
For every $\lambda>0$, positive semidefiniteness of $A_\sigma$ gives
\[
 \lambda(A_\sigma+\lambda I)^{-1}
 =I-(A_\sigma+\lambda I)^{-1}A_\sigma,
 \qquad
 \norm{(A_\sigma+\lambda I)^{-1}}_2\leq\lambda^{-1}.
\]
The first identity proves the limit and, for any class difference $v$, gives
\[
 \left|\lambda r(\sigma)^\tr(A_\sigma+\lambda I)^{-1}v
              -r(\sigma)^\tr v\right|
 \leq\frac{\norm{r(\sigma)}_2\norm{A_\sigma v}_2}{\lambda}.
\]
The stated threshold leaves a positive competing margin. Since an unbounded
penalty interval includes penalties above this finite threshold, the requested
whole-region agreement is impossible.
\end{proof}

Consequently, a necessary condition for whole-region agreement is
$k_a(\sigma,x)\geq k_b(\sigma,x)$ for every competitor and every
$0<\sigma\leq t$. A strict violation at one examined width suffices for
exclusion. An exact tie does not establish exclusion. Conversely, when class
$a$ strictly wins the limiting vote at a fixed width, the same error bound,
applied to each competing class, gives a finite penalty above which $a$
strictly wins every fitted decision at that width.

\subsection{Agreement on local width intervals and penalty tails}
\label{sec:local-tail}
The limiting vote also gives a positive certificate for local width intervals.
This extends the certified-region geometry to $[s,t]\times[\Lambda,\infty)$
with $s>0$. Both $r(\sigma)$ and $Z_\sigma$ are entrywise nonnegative and
nondecreasing in width. Their class scores therefore satisfy
$k_c(s,x)\leq k_c(\sigma,x)\leq k_c(t,x)$ for $s\leq\sigma\leq t$.
Define
\[
 m_{ab}(s,t,x)=k_a(s,x)-k_b(t,x).
\]
For each pair, enclose the signed training difference by
\[
 \ell_{ab}=Z_s[:,a]-Z_t[:,b],\qquad
 u_{ab}=Z_t[:,a]-Z_s[:,b].
\]
For a vector $z$, let $z_+=\max(z,0)$ and $z_-=\min(z,0)$ entrywise;
the negative part retains its sign. Set
\begin{align*}
 z^-_{ab}&=A_s(\ell_{ab})_++A_t(\ell_{ab})_-,\\
 z^+_{ab}&=A_t(u_{ab})_++A_s(u_{ab})_-,\\
 \rho_{ab}(s,t)&=
 \left\|\max\left(|z^-_{ab}|,|z^+_{ab}|\right)\right\|_2.
\end{align*}
All vector maxima and absolute values here are entrywise.

\begin{theorem}[Local-width, large-penalty agreement]
\label{thm:local-tail}
Suppose $0<s\leq t$ and $\Lambda>0$. If, for every class $b\ne a$,
\[
 m_{ab}(s,t,x)>
       \frac{\norm{r(t)}_2\rho_{ab}(s,t)}{\Lambda},
\]
then the unique winning class is $a$ for every
$(\sigma,\lambda)\in[s,t]\times[\Lambda,\infty)$ with finite $\lambda$.
\end{theorem}
\begin{proof}
The entrywise monotonicity of $H_\sigma$ gives
$A_s\leq A_\sigma\leq A_t$ entrywise. For
$v_\sigma=Z_\sigma[:,a]-Z_\sigma[:,b]$, we have
$\ell_{ab}\leq v_\sigma\leq u_{ab}$. Nonnegativity of $A_\sigma$
first gives $A_\sigma\ell_{ab}\leq A_\sigma v_\sigma
\leq A_\sigma u_{ab}$. Splitting each endpoint vector into its signed
positive and negative parts then yields
\[
 z^-_{ab}\leq A_\sigma v_\sigma\leq z^+_{ab},\qquad
 \norm{A_\sigma v_\sigma}_2\leq\rho_{ab}(s,t).
\]
The resolvent bound in Proposition~\ref{prop:ridge-limit}, the vote-margin
bound, and $\norm{r(\sigma)}_2\leq\norm{r(t)}_2$ imply
\[
 \lambda r(\sigma)^\tr
 (W_{\sigma,\lambda}[:,a]-W_{\sigma,\lambda}[:,b])
 \geq m_{ab}(s,t,x)-
       \frac{\norm{r(t)}_2\rho_{ab}(s,t)}{\Lambda}>0.
\]
Positive score scaling preserves the unique winning class.
\end{proof}

The certificate requires endpoint feature products and query distances.
Products with $A_s$ or $A_t$ can be computed as
$A_\sigma z=H_\sigma^\tr(H_\sigma z)$, so output fitting and inversion
are unnecessary. At $s=t$, the interval enclosure is exact:
$\rho_{ab}(t,t)=\norm{A_t(Z_t[:,a]-Z_t[:,b])}_2$. At every width
where the nearest-prototype label strictly wins the limiting vote, continuity
therefore provides a nonzero local width interval and a penalty tail with a finite floor
certified by Theorem~\ref{thm:local-tail}.

A simpler correction bound uses the positive class sum,
$\rho^+_{ab}(s,t)=\norm{A_t(Z_t[:,a]+Z_t[:,b])}_2$.
Nonnegativity and the endpoint bounds imply
$\rho_{ab}(s,t)\leq\rho^+_{ab}(s,t)$, so replacing $\rho_{ab}$ by
$\rho^+_{ab}$ gives a sufficient comparator on the same region.
When every $m_{ab}$ is positive, any penalty floor strictly exceeding
$\max_{b\ne a}\norm{r(t)}_2\rho_{ab}/m_{ab}$ passes the signed test.

\subsection{Certified, excluded, and unresolved requested regions}
\label{sec:partition}
For a fixed query and requested full-width region, the sufficient theorem
returns a certified case. A strict limiting reversal at any included width,
or a strict competing fitted score at any included finite parameter pair,
returns an excluded case: the region contains a disagreement and therefore
fails the uniform-agreement property. These two conclusions are mutually
exclusive. A query receiving neither conclusion is unresolved. Taking the
union of strict exclusions over examined widths preserves their validity
and makes the excluded set nondecreasing with the upper bandwidth.
The complement of the excluded set bounds attainable full-region agreement
coverage from above. The local-width theorem supplies a separate positive
certificate for its explicitly stated band and penalty tail.


\section{Evaluation design}
\subsection{Questions, datasets, and fixed sampling}
The frozen evaluation measures three properties: the fraction of queries certified over each declared region; the change in coverage between the three analytical bounds; and the relationship between certification, observed nearest-prototype agreement, and class accuracy. Numerical predictions at region endpoints and interior points check the implementation against Theorem~\ref{thm:agreement}. The proof establishes the continuous statement.

The datasets are MNIST handwritten digits \citep{lecun1998} and Fashion-MNIST apparel images \citep{xiao2017fashion}, each with ten classes and $28\times28$ grayscale inputs. Each seed uses 20,000 stratified training images and reserves 5,000 validation images. All pixels are converted to float64, divided by 255, and flattened consistently in MATLAB column order. Euclidean distances use the 784 resulting coordinates. The fixed seeds are 11, 22, 33, 44, and 55.

The MNIST training and validation partitions were retained from an earlier exploratory bandwidth study. Their union across seeds contains 55,973 of the 60,000 official training images. The 4,027 remaining images form a common evaluation set, with no overlap with any retained training or validation partition. The original 10,000-image MNIST test set is excluded from this prospective evaluation. Fashion-MNIST uses stratified training partitions from its official training set and the official 10,000-image test set for evaluation. In the frozen study, validation partitions were retained as disjoint unused partitions and no model selection was performed. Section~\ref{sec:r3-design} subsequently uses those partitions for a post hoc diagnostic of archived readouts. Canonical dataset payloads and local observation order were checked, and the archive retains hashes and index manifests.

Within each seed, the prototype budgets are $K\in\{50,100,200\}$. For MNIST, the earlier 100 randomly selected training prototypes are preserved; the first 50 form the smaller budget, and 100 further training rows are appended using seed $s+200000$. Fashion-MNIST draws 200 training rows using seed $s$ and uses prefixes for the smaller budgets. Prototype coordinates and labels remain fixed across all widths, penalties, and bound comparisons within each dataset--seed--budget group. Duplicate coordinates are recorded and retained. This creates 30 groups with matched comparisons between methods.

\subsection{Parameter regions and fitted references}
The prototype distance scale is
\begin{equation}
 s_C=\frac{\operatorname{median}_{j<k}\norm{c_j-c_k}_2}{\sqrt{2}}.
 \label{eq:scale}
\end{equation}
The seven bandwidth endpoints are the absolute value $t=1$ and $t/s_C\in\{1/16,\allowbreak1/8,\allowbreak1/4,\allowbreak1/2,\allowbreak1,\allowbreak2\}$. Each combines with $\lambda_0\in\{0,1,100\}$ to define a complete region. Thus the experiment contains 630 regions and 1,890 method-specific coverage records. Every declared condition is retained.

For implementation checks, each region is probed at bandwidths $\{t/2,3t/4,t\}$. The penalty probes are $\{0,1,100\}$ when $\lambda_0=0$, $\{1,10,100\}$ when $\lambda_0=1$, and $\{100,1000,10000\}$ when $\lambda_0=100$. The Cartesian product gives nine checks per certified query and region. Fitted references use thin singular value decomposition (SVD). Positive penalties apply the filter $s/(s^2+\lambda)$ to each singular value $s$; zero penalty uses a pseudoinverse with cutoff $\max(N,K)\epsilon s_{\max}$, where $\epsilon$ is float64 machine precision. Numerical ranks are retained. Gaussian training features are unscaled. Query features use the distance-gap normalization in Eq.~\eqref{eq:query}, preserving decisions and preventing a common nearest-response factor from underflowing.

The implemented certificate requires each margin gap to exceed $10^{-10}(1+|L_b|+|n_tD_{ab}|)$. It also requires the difference between the smallest two squared distances to exceed $10^{-12}\max(1,d_{(1)}^2,d_{(2)}^2)$. Otherwise it abstains. These rules are shared by all three methods. The code checks bound ordering and inclusion of certificate masks and stores the minimum raw class-pair margin for every query-region pair.

\subsection{Measures and execution record}
For a group with $M$ evaluation queries, coverage is $M^{-1}\sum_{i=1}^M I_i$, where $I_i$ indicates certification. Applicability records whether $\gamma_t+\lambda_0>0$. Observed agreement is the fraction of queries for which a fitted readout and the nearest prototype return the same label; accuracy compares a fitted label with the held-out target. Coverage among observed agreements is reported with the agreeing queries as denominator. These measures distinguish a guarantee obtained from the sufficient condition, direct agreement at a fitted point, and correctness against the dataset label.

The protocol and implementation hashes were frozen on 6 November 2026 before evaluation. Predictions, weights, class-pair bounds, and certificate masks for all 30 groups were generated and hashed before evaluation targets were used for scoring. The loaded MNIST label array supplied training labels; entries corresponding to the reserved evaluation indices were excluded from fitting, selection, and scoring until that prediction freeze. All comparisons preserve the same data, prototypes, and parameter region.

Means and sample standard deviations summarize the five seed outcomes for each fixed condition. Comparisons between methods are paired within a group. Query-region and query-probe totals count repeated uses of the same images and therefore serve as implementation-check denominators. The archived run used Python 3.12.14, NumPy 2.3.5, Linux x86-64, and single-thread limits around the computational loop. Timing records separate shared distances, certificate preparation with query evaluation, and readout factorization.

\subsection{Supplementary bandwidth-frontier analysis}
The original grid places its positive scaled endpoint at $s_C/16$ and its next endpoint at $s_C/8$. To locate the transition, a post hoc analysis on 7 November 2026 retains all 30 groups, their original indices and prototypes, and all three penalty floors. Its protocol and code were frozen before the supplementary run. It evaluates 67 geometrically spaced widths from $s_C/32$ to $s_C/4$, with 22 intervals per doubling and adjacent ratio $2^{1/22}\simeq1.03201$. The original widths $s_C/16$, $s_C/8$, and $s_C/4$ are included exactly.

For each query, floor, and method, the stored frontier bracket contains the largest passing grid width and the next failing width. Monotonicity of the class-pair bounds and query margins makes acceptance downward closed. The numerical guard preserves this ordering: for a passing comparison, $L>A\geq0$, where $A=n_tD_{ab}$, and the guard is equivalent to $(1-\epsilon)L>(1+\epsilon)A+\epsilon$ with $\epsilon=10^{-10}$. The left side decreases with width and the right side increases. Numerical mask nesting is also checked. Queries already failing at the first width are recorded as left-censored; those passing the final width are recorded as right-censored. Each observed transition is bracketed with an upper-to-lower width ratio no greater than 1.0321.

For each group and floor, sharp masks select three operating points: the largest grid width retaining any certification, at least 50\% coverage, and at least 90\% coverage. Selection uses no evaluation targets. A readout is fitted at each selected width and $\lambda=\lambda_0$; all predictions are frozen before scoring. Accuracy and nearest-prototype agreement use all queries, with certified-subset accuracy retained separately. The comparison reference is the highest observed accuracy among the 16 distinct widths already probed in the frozen study at the same penalty. This reference describes the sampled accuracy range. All five seed results and their paired differences are retained.

\input{new_methods.tex}
\input{r3_methods.tex}
\input{r4_methods.tex}

\section{Results}
\input{results.tex}
\input{frontier_results.tex}
\input{new_results.tex}
\input{r3_results.tex}
\input{r4_results.tex}

\section{Discussion}
The combined analysis distinguishes a nearest-prototype summary that is guaranteed across a family from one that the family itself contradicts. The residual certificate controls all widths below an endpoint and every penalty above a floor. The large-penalty obstruction shows when some finite member of that same family must choose a competing class. Their combination adds information to a sufficient certificate's abstention: excluded regions contain a demonstrated disagreement, whereas unresolved regions await a decisive test.

The region diagnostic separates two empirical phenomena. Immediately beyond the certification frontier, almost the entire query set remains unresolved. At broader widths, strict witnesses exclude most requested full-width regions. These observations establish that the first loss of sufficient coverage and widespread identified failure of uniform agreement occur in different bandwidth ranges. A tighter sufficient bound could resolve some of the remaining cases; any sound bound must exclude queries with a strict disagreement witness from its accepted set.

The large-penalty limit explains a source of incompatibility. The scaled readout approaches $Z_\sigma=H_\sigma^\tr Y$, and the query votes $r(\sigma)^\tr Z_\sigma$ aggregate labeled observations through their prototype responses. Non-prototype observations therefore influence the limiting decision even though the reference label belongs to one selected prototype. Increasing regularization can expose a pooled class vote that opposes that reference. The finite threshold in Proposition~\ref{prop:ridge-limit} makes this mechanism a statement about actual finite classifiers. The validation comparison places these sufficient bounds far beyond the competitive sampled penalties. The complete archived paths exhibit both earlier finite disagreements and later returns to agreement, including five alternating signs across six penalties. This distinguishes a guaranteed disagreement tail from the observed sequence of finite decisions. \ref{app:exact-ridge-examples} supplies exact constructed cases for strict reversal and tied-limit behavior.

The local-width theorem uses the same limiting mechanism constructively. Positive vote margins across a bandwidth band, together with a signed enclosure of the resolvent correction, certify a penalty tail with a finite floor. The measured coverage at broader widths establishes that this region geometry can certify additional nearest-label behavior. The matched-coverage experiment places its selected-set accuracy above random selection and ordinary score margins. Agreement-gated distance selection is more accurate in three of the four comparisons. The MNIST endpoint at $t=s_C$ and $\lambda=10\norm{A_t}_2$ has a constant class-1 readout, and its certified set equals the entire endpoint-agreement pool. The class-specific reconstruction makes this composition explicit: the MNIST certified subset at that endpoint has 10\% selected-class macro accuracy across ten represented true classes. On Fashion-MNIST, the stronger selector distributes correct accepted predictions across more classes and selects substantially different queries at equal coverage. These measurements separate continuous persistence, selected-subset accuracy, and class-balanced retention.

Exact overlap analysis further identifies the source of incremental exclusions. Near $t=s_C$, the original limiting witnesses add many cases to the original finite probes, and almost all of those additions are covered once the archived local-tail models are included. At $2s_C$, limiting witnesses retain a larger unique set on Fashion-MNIST, where the archive has fewer high-penalty finite models. Complementarity is therefore a measurable relation between the analytical witness and a specified finite sampling design. The limiting calculation provides a training-derived route to the exclusion; finite readouts expose many of the same disagreements. The fixed replay order sharpens this comparison: expanded high-penalty probes recover at least 95\% of the original limiting-witness cases within mean 15.6 and 20.2 checks at $t=s_C$, whereas the original finite probes fail to reach that target in every seed.

The original bound comparisons explain which information matters within the full-width construction. Class-wise overlap retains the distribution of training responses by label, whereas the Frobenius comparator compresses them into one norm and class counts. The larger certified sets establish the effect of this retained structure under identical data and query margins. The exact two-endpoint envelope further improves the analytical bound. Its masks match the basic class-wise masks in the frozen grid, and the denser grid resolves small additional sets near some transitions. The larger empirical change comes from class-wise information; the envelope removes a continuous penalty maximization and supplies a further local gain. Even at the widest original passing endpoint, the Gaussian response at median prototype separation is $2.96\times10^{-32}$, placing the full-width sufficient construction in a near-isolated-prototype regime.

The validation-selected analysis supplies a direct empirical answer about competitive operation. Every model retained by each of the fixed 90\%, 95\%, and 99\% rules has zero sharp full-width coverage. The selected models improve accuracy and cross-seed prediction stability as prototype budget increases, while the full-width certificate remains empty. Thus improved finite-model prediction and certified family persistence follow distinct empirical patterns in the archived grid. The broad region exclusions additionally show that many competitive endpoints belong to families containing a strict nearest-label disagreement.

These results connect established RBF and parameter-path analyses through a specific decision-region object. Coefficient-regularized RBF theory addresses statistical behavior \citep{que2016}; narrow-bandwidth kernel arguments explain nearest-neighbour dominance \citep{radhakrishnan2023}; prototype architectures and learned Gaussian geometry address interpretation through model structure \citep{saralajew2025,dagostino2024}. Regularization-path guarantees control solution or validation behavior as penalties vary \citep{shibagaki2015,ndiaye2019}. Here the varied feature matrix and penalty are analyzed together to establish nearest-prototype label persistence, strict incompatibility, or unresolved status before output fitting. The signed local certificate adds a complementary region construction using the same training geometry.

\section{Conclusion}
Training-derived quantities can both certify and exclude nearest-prototype agreement in Gaussian ridge RBF families. A class-wise residual, spectral floor, and exact penalty envelope certify complete regions $(0,t]\times[\lambda_0,\infty)$. A strict competing vote in the large-penalty limit excludes such a region and supplies a finite disagreement threshold. A signed interval bound additionally certifies local bands $[s,t]\times[\Lambda,\infty)$. The two-dataset evidence separates the early loss of sufficient certification, broad-width incompatibility, and positive local-tail coverage. All 328 validation-competitive models under the primary rule have zero full-width coverage; the finite paths exhibit earlier disagreements and subsequent returns to agreement. Exact witness unions and fixed-order replay quantify the relation between analytical exclusions and finite checks. Class-balanced matched-selection diagnostics explain the observed subset accuracy. The contribution is a computable analysis of continuous decision persistence with certified, excluded, and unresolved outputs, accompanied by an exhaustive assessment of its reach in the archived model families. Tighter signed bounds and alternative region geometries provide concrete routes for studying the unresolved frontier.

\section{Limitations}
The guarantees concern equality with a nearest-prototype label. Class correctness is evaluated separately against held-out targets; human understanding and causal explanatory value were not measured. The model uses fixed labeled training prototypes, one shared Euclidean Gaussian width, the sum-of-squares ridge objective, and no intercept. Learned features, moving prototypes, class-dependent penalties, anisotropic kernels, and input perturbations require corresponding derivations. The analytical complexity statement concerns the proposed accumulation, whereas the experimental implementation retains dense feature matrices and supplies descriptive timing rather than a streaming-memory or speedup benchmark.

The full-width sufficient condition discards information through absolute class-pair majorization, a prototype-block eigenvalue floor, removal of additional same-class query responses, and norm inequalities. The diagnostic examines finitely many limiting widths and archived finite readouts; unresolved queries can contain disagreements at other parameter values. Limiting ties receive no exclusion from the strict test. The local-tail certificate addresses $[s,t]$, and the full-width certificate addresses $(0,t]$; their coverage refers to these different families. Large-penalty local coverage and certified-subset accuracy do not measure all-query predictive performance.

Two ten-class image datasets and randomly selected prototype subsets determine the empirical scope. The post hoc diagnostic, local-tail, reversal-bound, selector, and ordered-probe analyses use the ten $K=100$ groups. Finite-path, validation-competitive, and budget-stability analyses use all 30 original groups. The 90/95/99\% rules describe relative validation accuracy within the archived grid. Sampled first-detection brackets do not identify unique or earliest continuous crossings, and all-sample agreement does not certify unsampled penalties. The sufficient reversal bound supplies an upper threshold for guaranteed disagreement. Ordered-probe prefix counts measure retrospective set recovery; they supply neither a prospective stopping guarantee nor a measured computational speedup. The selector comparisons describe the fixed endpoint models and sample. No predictive or model-selection improvement is established by the certificates. Seeds share evaluation images and have overlapping training pools, so sample standard deviations describe variation across those realizations. Repeated query-region and query-probe checks reuse the same images. Float64 arithmetic and positive guards were used without formally outward-rounded interval arithmetic. The finite probes check implementation behavior, and the proofs establish the continuous statements. The supplementary frontier brackets concern the guarded numerical certificate; the sampled accuracy maximum is a descriptive reference selected using evaluation targets. The present artifact verification reconstructs archived measurements, while the earlier complete-pipeline rerun is documented in the frozen archive.

\appendix
\section{Duplicate inputs and the small-width reference}
\label{app:duplicates}
The reference $W_0(\lambda)$ is also the small-width limit under a specific separation condition. If every selected prototype coordinate occurs only at its selected training row, then $H_\sigma\to S$ as $\sigma\to0$. For fixed finite $\lambda\geq0$, continuity gives $W_{\sigma,\lambda}\to Y_C/(1+\lambda)$. A query with a unique nearest prototype has $r(\sigma)\to e_{j_0}$, so its limiting class is the prototype's label. The finite-region theorem uses the actual overlap terms and its strict condition directly.

A constructed example shows why non-prototype training observations matter. Let training coordinates be $(0,0,0,3)$ with labels $(0,1,1,1)$, and select the first and fourth rows as prototypes. The first prototype has label zero, and two further observations at the same coordinate have label one. As width tends to zero, the readout at penalty one converges to
\[
 W\longrightarrow
 \begin{pmatrix}1/4&1/2\\0&1/2\end{pmatrix}.
\]
A query at zero consequently receives class one although its nearest selected prototype has class zero. Here $P_t$ retains the two coincident non-prototype observations, and the certificate abstains. The example explains the role of those training-label contributions; it is a constructed algebraic case.

\section{Synthetic verification of the implementation}
Two separately written numerical checks tested the derivation before the image evaluation. The first generated 100 small training datasets and evaluated 6,555 query-region cases, including 4,394 certified regions and 163,875 fitted comparisons. It found zero certified-label violations, residual-majorant failures, spectral-bound failures, or sharp-over-basic bound violations at its numerical tolerance. A separate implementation generated 60 further geometries and checked 117,975 certified query-parameter combinations with zero prediction disagreements. Its maximum residual-identity discrepancy was $1.42\times10^{-14}$; all 6,976 finite directed class-pair comparisons satisfied the sharp/basic ordering, with 3,290 strict improvements. The duplicate-input example and nearest-distance tie handling were also exercised. The archived scripts preserve these generated cases and checks separately from the two-dataset evaluation.

An additional independently written check of Proposition~\ref{prop:ridge-limit} and Theorem~\ref{thm:local-tail} generated 250 Gaussian training geometries with 24 observations, six selected prototypes, three classes, and 12 queries per geometry. It verified 15,750 signed interval enclosures, 3,000 resolvent inequalities, 2,930 strict-limit finite reversals, and 33,660 fitted predictions in certified local bands, with zero failures. The exact Gaussian examples in \ref{app:exact-ridge-examples} separately exercise a strict obstruction and a tied limit. The revision supplement retains the independent verification script and its complete output.

\input{exact_ridge_examples.tex}

\section*{Funding}
This research did not receive any specific grant from funding agencies in the public, commercial, or not-for-profit sectors.

\section*{Declaration of competing interest}
The author declares no known competing financial interests or personal relationships that could have appeared to influence the work reported in this paper.

\section*{CRediT authorship contribution statement}
Meshaal Mouawad: Conceptualization, Methodology, Software, Validation, Formal analysis, Investigation, Data curation, Visualization, Writing -- original draft, Writing -- review and editing.

\section*{Data and code availability}
The code, indexed data splits, prototype identities, frozen predictions, certificate masks, fitted weights, supplementary bandwidth-frontier records, dependency specifications, and verification scripts are provided in the accompanying Supplementary Research Artifacts archive. Dataset provenance and licenses are retained in the archive. The original MNIST and Fashion-MNIST sources are identified in the evaluation design and archive documentation.

The additional diagnostic and local-tail protocols, code, per-query records, fitted readouts, and independent synthetic checks are provided in the accompanying revision supplement. It also includes the follow-up validation predictions, sufficient-reversal distributions, exact witness-overlap masks, matched-cardinality selection records, class distributions, and independent numerical reconstruction. The all-budget extension adds validation-selected and competitive-set records for every original group, per-class composition, complete signed finite-path margins, fixed-order exclusion increments, paired budget differences, cross-seed overlaps, plots, and independent reconstruction scripts. The frozen original archive and the additional supplement retain separate provenance records.

\section*{Declaration of generative AI and AI-assisted technologies in the manuscript preparation process}
During preparation of this work, the author used Elsevier LeapSpace, other tools to support literature discovery, manuscript organization, language refinement, and consistency checking. The authors take responsibility for the content of the article.

\bibliographystyle{elsarticle-harv}
\begingroup\raggedright
\setlength{\bibsep}{2pt}
\bibliography{references}
\endgroup
\end{document}
